% ============================================================
\section{Protocol Specification}
\label{sec:protocol}

This section specifies M\textsuperscript{2} as a deployable system: four immutable contracts, a genesis ceremony, the operations they expose, and a threat model that delimits the conditions under which the formal guarantees of Section~\ref{sec:formal-invariants} hold. Notation follows Appendix~\ref{app:symbols}.

\paragraph{System.} M\textsuperscript{2} is composed of four contracts---a token, a treasury, an asymmetric-fee LP hook, and a revenue router---together with one off-protocol \emph{operator} who routes business revenue (denominated in the backing stable) into the router.\footnote{The Solidity reference implementation is at the Zenodo artifact referenced in the abstract; the bytecode is the contract.} Each contract is deployed once with no admin key, governance hook, upgrade path, or parameter-change interface. Figure~\ref{fig:architecture} shows the system. Dynamics: (i) the operator deposits revenue $X$ into the router, which atomically deposits $X/2$ into the treasury and uses $X/2$ for a buy-and-burn on the protocol-owned LP; (ii) holders invoke either \texttt{redeem} on the token (atomic, no cooldown, no fee) or an LP swap (asymmetric per-direction fee enforced by the hook); (iii) any external caller may invoke \texttt{collectFees}, which routes $99.75\%$ of token-side fees to burn and $99.75\%$ of stable-side fees to the treasury, with $0.25\%$ per side paid as a caller bounty.

\begin{figure}[t]
\centering
\resizebox{0.6\textwidth}{!}{%
\begin{tikzpicture}[
  node distance=1.5cm and 2.0cm,
  contract/.style={draw, thick, rounded corners=2pt, minimum width=3.2cm, minimum height=0.9cm, align=center, font=\small},
  business/.style={draw, thick, dashed, rounded corners=2pt, minimum width=2.4cm, minimum height=0.6cm, align=center, font=\small},
  every node/.style={font=\small},
  flow/.style={->, thick, >=stealth},
  feeflow/.style={->, thick, >=stealth, dashed, gray!70!black},
  immut/.style={fill=blue!5},
]
  \node[contract, immut] (token) {Token (ERC-20) \\ \scriptsize fixed supply $\Sup_0=10^9$, burn, redeem};
  \node[contract, immut, right=of token] (treas) {Treasury \\ \scriptsize stable reserve $\Tres$, one-way in};
  \node[contract, immut, below=of token] (hook) {LP + V4 Hook \\ \scriptsize asym.\ fees, anti-MEV};
  \node[contract, immut, below=of treas] (router) {Revenue Router \\ \scriptsize 50/50 split, immutable};
  \node[business, below=0.9cm of router] (biz) {Business};
  \draw[flow] (token) -- node[above, font=\scriptsize]{redeem} (treas);
  \draw[flow] (hook) -- node[left, font=\scriptsize]{burn} (token);
  \draw[flow] (router) -- node[above, font=\scriptsize]{50\%} (treas);
  \draw[flow] (router) -- node[above, font=\scriptsize]{50\%, buy} (hook);
  \draw[flow] (biz) -- node[right, font=\scriptsize]{routeRevenue} (router);
  \draw[feeflow] (hook.west) to[out=180, in=200, looseness=1.5]
      node[left, font=\scriptsize, align=right]{collectFees:\\token burn} (token.west);
  \draw[feeflow] (hook.south east) to[out=-20, in=-90, looseness=1.4]
      node[below right=2pt, font=\scriptsize, align=left]{collectFees:\\stable to $\Tres$} (treas.south east);
\end{tikzpicture}}
\caption{System architecture. Blue-tinted boxes are immutable contracts; the dashed box is the off-protocol operator. Solid arrows mark primary flows; dashed gray arrows mark the permissionless \texttt{collectFees} realization paths (token side: $99.75\%$ burned, $0.25\%$ to caller; stable side: $99.75\%$ to treasury, $0.25\%$ to caller). All inflows to Treasury are one-way: stablecoin can leave only via \texttt{redeem}, which simultaneously burns the redeemed tokens.}
\label{fig:architecture}
\end{figure}

\paragraph{Token, redemption, floor precision.} The token is an ERC-20 with $18$ decimals and total supply $\Sup_0 = 10^9$ minted to the deployer in a single transaction; the constructor-set genesis-mint sentinel blocks all further minting. Three callers can \texttt{burn}: the hook (token-side LP fees), the router (buy-and-burn), and the token's own \texttt{redeem}. \texttt{redeem(uint256 amount)} is atomic, non-pausable, with no cooldown, fee, or maximum: it burns the caller's tokens, computes the payout as $\texttt{amount} \times \texttt{floorPrice}()$, and instructs the treasury to transfer via the privileged \texttt{payRedemption} interface. The floor returns an 18-decimal fixed-point value regardless of the stable's native decimals $\dstab$:
\begin{equation}
\texttt{floorPrice}() = \frac{\Tres \cdot 10^{18 + 18 - \dstab}}{\Sup},
\label{eq:floor-precision}
\end{equation}
with the wide intermediate in \texttt{uint256}; $10^{42} \ll 2^{256}$ rules out overflow. Integer truncation can leave at most $\Sup - 1$ smallest units of dust in the treasury's favor per call; Lemma~\ref{lem:integer-redemption} shows the residual strictly raises the floor. The treasury exposes exactly one privileged entry point, \texttt{payRedemption}, callable only by the token contract; no admin, withdraw, pause, upgrade, or governance hook. Stablecoin leaves only through \texttt{payRedemption} in lockstep with a token burn.

\paragraph{Asymmetric-fee hook and V2-equivalent algebra.} The protocol-owned LP is a Uniswap v4 pool initialized with \texttt{DYNAMIC\_FEE\_FLAG} and a single \emph{full-range} position spanning $[\text{MIN\_TICK}, \text{MAX\_TICK}]$. Under this configuration v4's active liquidity equals the position's at every reachable price and per-tick concentrated-liquidity bookkeeping never engages: $(\Lt, \Ls)$ are V2-style reserves of a single piecewise-constant liquidity segment, and $\kInv = \Lt \cdot \Ls$ is exact modulo the asymmetric-fee deduction. The position-owning contract has no removal interface, so liquidity is provably unwithdrawable. On every swap, the hook's \texttt{beforeSwap} returns a fee override determined by direction:
\begin{equation}
\text{fee}(\text{direction}) =
\begin{cases}
0.10\% & \text{if input is stable (buy)} \\
3.00\% & \text{if input is token (sell)}
\end{cases}
\label{eq:fee-asymmetry}
\end{equation}
The fee is subtracted from input before the curve and folded into the input-side accumulator, denoted $(\Phit, \Phis)$ in Section~\ref{sec:formal-invariants}.\footnote{V4 stores the accumulator as Q128.128 \texttt{feeGrowthGlobalX128}; the raw-token framing here is dimensionally equivalent.} The $3.00\%$ sell fee converts every adversarial LP exit into common-pool value. A holder's two exit paths are both floor-non-decreasing: \texttt{redeem} preserves the floor exactly (Case~3), and the LP-sell path holds it flat at swap and strictly raises it at the next \texttt{collectFees} (Cases~5,~7). Protocol-initiated buys use private-orderflow routing, a uniform delay $\tau \sim U[0, 12\,\text{h}]$, and optional TWAP splitting; residual MEV leakage is bounded by Theorem~\ref{thm:mev-bound} (Appendix~\ref{app:auxiliary}).

\paragraph{Revenue router and \texttt{collectFees} distribution.} The router exposes a single function callable only by the immutably-set depositor: it pulls a stable amount, deposits $\lfloor \cdot/2 \rfloor$ into the treasury, and uses $\lceil \cdot/2 \rceil$ for a buy-and-burn via the hook in the same transaction. The split is a bytecode constant---no setter, upgrade path, or governance hook. The hook exposes a permissionless \texttt{collectFees}: it acquires the pool-manager lock, calls \texttt{modifyLiquidity} with zero delta, obtains $(\Kreal, \Ureal) = (\Phit, \Phis)$, and distributes:
\begin{align}
\text{caller bounty (token):} \quad & K_b = \lfloor 0.0025 \cdot \Kreal \rfloor, \\
\text{caller bounty (stable):} \quad & U_b = \lfloor 0.0025 \cdot \Ureal \rfloor, \\
\text{token burn:} \quad & K_{\text{burn}} = \Kreal - K_b, \\
\text{stable to treasury:} \quad & U_{\text{treas}} = \Ureal - U_b.
\end{align}
\texttt{collectFees} does \emph{not} route through the revenue router: its destinations are unambiguously floor-positive and the $50/50$ split applies only to direct business revenue. The $0.25\%$-per-side bounty pins the equilibrium inter-call interval at Theorem~\ref{thm:collect-equilibrium}'s break-even point.

\paragraph{Genesis ceremony and threat model.} The genesis transaction atomically deploys the four contracts, mints $\Sup_0 = 10^9$ tokens, deposits $\Tres_0$ into the treasury, creates the v4 pool with a single full-range position holding $L_{t,0} = 7.5 \cdot 10^8$ tokens and $L_{s,0}$ stable, and transfers the remaining $25\%$ to a vesting vault. The deployer-chosen seeds are constrained by
\begin{equation}
\frac{\Tres_0}{\Sup_0} \;=\; \frac{L_{s,0}}{L_{t,0}},
\label{eq:genesis-constraint}
\end{equation}
forcing $\Floor_0 = \Spot_0$ at launch. The canonical instance uses $\Tres_0 = \$1$M and $L_{s,0} = \$750$k, giving $\Floor_0 = \Spot_0 = \$0.001$/token, with design constants $f_b = 0.10\%$, $f_s = 3.00\%$, $0.25\%$-per-side bounty, $50/50$ router split, $75/25$ LP/vesting split. Table~\ref{tab:threat-model-condensed} delimits which adversaries are in scope; the full $14$-row surface is in the companion supplementary file \texttt{supplementary-threat-model.tex} bundled with the Zenodo artifact.

\begin{table}[t]
\centering
\footnotesize
\caption{Body threat model (six load-bearing rows). ``In scope'' means the protocol's invariants are claimed to hold against the indicated capability; ``out of scope'' acknowledges residual risk requiring separate analysis. The full $14$-row surface is in the companion file \texttt{supplementary-threat-model.tex}.}
\label{tab:threat-model-condensed}
\begin{tabular}{p{2.8cm} p{4.8cm} p{2.0cm} p{3.4cm}}
\toprule
\textbf{Adversary} & \textbf{Capability} & \textbf{Scope} & \textbf{Where treated} \\
\midrule
Rational holder & Any composition of \texttt{redeem}, LP swap, transfer; observe public state & In scope & Thm~\ref{thm:floor-monotonicity}, Thm~\ref{thm:redemption-fairness} \\
Coordinated coalition (any size up to total supply) & Submit any sequence of in-protocol operations; coordinate timing & In scope (simultaneous-strategy class) & Thm~\ref{thm:bank-run}; sequenced class is Section~\ref{sec:limitations} (P3) \\
MEV searcher (single block) & Reorder, sandwich, frontrun within a block; pay priority gas & In scope (bounded leakage) & Thm~\ref{thm:mev-bound}; hook mitigations \\
Operator (revenue depositor) & Choose timing and magnitude of \texttt{routeRevenue}; choose not to call & Partially in scope: floor preserved at zero revenue; growth operator-dependent & Section~\ref{sec:economic-analysis}; Section~\ref{sec:limitations} \\
Backing-stable issuer & Freeze treasury balance; depeg the stable; halt transfers & Out of scope (exogenous) & Section~\ref{sec:limitations}; USDC March-2023 depeg as canonical event \\
Smart-contract bug + L1/L2 execution layer & Violate any invariant via implementation defect; halt chain; long-term censorship & Out of scope (no post-deployment fix; substrate assumption) & Section~\ref{sec:limitations}; audit + mechanized-verification path \\
\bottomrule
\end{tabular}
\end{table}

M\textsuperscript{2} is trust-minimized, not trustless: issuer freeze, smart-contract bugs, multi-block MEV under builder collusion, operator fraud, cross-block \texttt{collectFees}-adjacent MEV, and third-party oracle consumers are residual surfaces catalogued in Section~\ref{sec:limitations}.
